\documentclass[11pt]{article}
\usepackage{amsmath,amssymb}
\usepackage[hidelinks]{hyperref}
\title{What misaligned factory logs can certify before dispatch}
\date{}
\begin{document}
\maketitle

A log can support two different conclusions: some authorised copy
physically escaped the interruption, and a particular actor could
identify an authorised copy before committing. Neither conclusion
implies the other just from timestamps. A small trace family gives a
sharp separation. Shared clock calibration certifies that some copy
survives on every history, but neither copy alone is guaranteed.
The first cutoff that guarantees a timely identifying receipt is
exactly two ticks. Treating the clocks as independent loses even
the physical-survival conclusion, without changing the individual
route readings.

\section{Measured source and limits}
The \href{https://zenodo.org/records/4034867}{public CONTEXT v2 release}
contains measured OPC-UA change histories and independent sensing
streams from a relay-production line. The ZIP directory shows four
runs: ground truth, shuttle dropout, missing parts, and pressure
loss, with six CSV files per run. \href{https://www.sciencedirect.com/science/article/pii/S1877050921003148/pdf}{Kaupp
et al.}, sections 3--4, describe separate acquisition clocks,
missing or out-of-order events, and a proposed approximate alignment
using first activity and shuttle travel. The approximately 30-second
travel time is not a certified upper or lower bound for an event
pair. The physical station sensor is not by itself a unique material
identifier, and an OPC-UA archival timestamp does not establish
when the deciding actor received the data. A workpiece code used
to choose an inspection routine only identifies the workpiece for
this analysis if a specific downstream linkage can be established.
No numerical CSV record or intervention in this release has been
selected or scored here. The release motivates the measurement
problem; it does not validate the synthetic timing example below.

For an event-level evaluation, record the origin clock or time
mark for each event, uncertainties with their source justification,
and the difference between an attempted effect, actual custody,
sensor read, log admission and receipt by an authorised actor.
Use bounded clock drift over an admitted episode if finite, and
separate captures that can be absent from captures that are
guaranteed but late. Where a proposed cross-clock match can be
wrong, admit the alternate matches as separate possible models
or withhold the corresponding conclusion. If two groups of clock
variables have no relative bounds, an arbitrary common shift within
either group preserves its local readings. In that model no
cross-group precedence follows from the logs alone. A point
alignment always imposes additional assumptions.

\section{Checkable chronology}
Use time variables \(t_v\), including local clock offsets or a
reference epoch when required. The source supplies bounds
\(t_v-t_u\leq b\) or \(t_v-t_u<b\); a fixed logged offset with
bounded delay can be encoded by two such bounds. Bounds within
and between clocks are kept jointly. The source premises must
have at least one model, so inconsistency is not accepted as
evidence of safe action.

For each acceptable route \(r\), let \(Q_r\) be conjunctions of
the applicable timing constraints; for example actual collection
before physical destruction, and an identifying report to the
deciding actor by its commitment. To establish
\(\forall t\in{\cal T},\,\exists r:Q_r(t)\), choose one failed
constraint per route and test all such tuples of complements on
the \emph{same} source model \({\cal T}\). This is exact because a
history fails all routes iff it clears one such tuple of tests.
The joint failure may be impossible even though each failure is
separately reachable. Each tuple is a simple difference-bound
feasibility problem. An explicit schedule is a counterexample;
negative cycles in all augmented systems certify success. In the
integer-clock model the complement of
\(t_v-t_u\leq b\) is \(t_u-t_v\leq-b-1\).
The independent checker also verifies a base-model witness.

For rational or continuous clock times, one must not replace
strict order by a unit-tick gap. The complement reverses the edge
and swaps strictness. If bounds have common denominator \(D\)
and there are \(n\) vertices, a strict bound can be decreased
by exactly \(1/(D(n+1))\) in the feasibility check. Indeed a
nonnegative simple cycle with positive original total has at
least \(1/D\) of slack and loses at most \(n/(D(n+1))\); a zero
cycle with any strict edge is infeasible regardless. This
reduction and an ordinary potentials solver yield an exact
rational witness without an arbitrarily chosen tolerance. The
statement is about satisfiability of these supplied difference
bounds, not about all clock drift, or about how to infer good
bounds from a site. Checking every failure tuple can be
exponential in the number of multi-constraint routes.

\section{Cut and fixed-route separation}
In an illustrative two-copy civil recovery model there are two
authorised preparations of the same obligation. The physical
clock offset \(x\in[-1,1]\) is shared by two station events.
On A the attempted pickup is at \(x\), while the fault is at
zero. On B pickup is attempted at zero, while the fault is at
\(x+1\). Capture strictly before the local fault is required;
tie means destruction first. Thus A succeeds when \(x<0\)
and B when \(x>-1\). At least one capture always succeeds,
but neither fixed route is safe on every history. Replacing
the common \(x\) by detached offsets \(x_A,x_B\) with the same
marginal interval admits simultaneous failure, e.g. \(0,-1\).
This is an ablation of physically justified clock dependence, not
an equally informed rival planner.

A positive report is generated only for an actual secure capture;
the source time bounds for an attempted pickup do not themselves
guarantee a positive message. When present, the authorised actor
gets the message after one to two ticks from the relevant pickup.
At cutoff one the history \(x=0\) and the B delay two has no timely
positive evidence; at cutoff two the actor always learns the
identity of at least one captured copy. Waiting until then is
only an admissible repair when the plant still allows that cut.
If only one input can reach the actor, even cutoff two cannot
guarantee such evidence. The controller use of a timely positive
report is explicit: choose an authorised route actually verified
at the cut. It neither chooses the logically existing route
without that evidence nor sends a copy prematurely. If a
locally trusted fail-safe selector can observe both effects, it
is a different, charged observation interface.

\section{Status}
\href{https://www.sciencedirect.com/science/article/pii/0004370291900066}{Dechter,
Meiri and Pearl} provide the standard temporal-network machinery.
A matched feasibility comparator ties. The result is a calibration
and audit method, not a new temporal solver or a general dynamic
controllability guarantee. Four Lean statements confirm the
fixture's exact rational cutoff and loss under detaching the
source offsets. The Python replays compare integer and rational
systems with exhaustive enumeration and verify explicit
failure schedules or negative cycles in the fixed cases.
This is development work for I1, not an independent compound
scenario or coverage disposition.
\end{document}
